\documentclass[11pt]{article}

\usepackage[a4paper,margin=1in]{geometry}
\usepackage{amsmath,amssymb,amsthm,mathtools}
\usepackage{microtype}
\usepackage[hidelinks]{hyperref}

\newcommand{\HH}{\mathcal H}
\newcommand{\E}{\mathbb E}
\newcommand{\Prob}{\mathbb P}
\newcommand{\R}{\mathbb R}
\newcommand{\inner}[2]{\left\langle #1,#2\right\rangle_{\HH}}
\newcommand{\norm}[1]{\left\lVert #1\right\rVert_{\HH}}
\newcommand{\pos}[1]{\left[#1\right]_+}

\DeclareMathOperator*{\argmin}{arg\,min}

\newtheorem{assumption}{Assumption}
\newtheorem{theorem}{Theorem}
\newtheorem{proposition}{Proposition}
\newtheorem{remark}{Remark}

\title{Finite-Sample Guarantees for Regularized RKHS
Pearson Classification}
\author{}
\date{}

\begin{document}
\maketitle

\section{Setting and classification score}

Let $\HH$ be a real separable reproducing kernel Hilbert
space with measurable feature map $\phi(x)=k(x,\cdot)$.

\begin{assumption}[Kernel and sampling]
\label{ass:setting}
The positive-semidefinite kernel satisfies
\[
    k(x,y)\ge 0,
    \qquad
    k(x,x)\le K
\]
for all $x,y$, where $K>0$.

For a fixed reference regime $P_i$ and an incoming
distribution $Q$, let
\[
    X_1,\ldots,X_n \overset{\mathrm{iid}}{\sim} P_i,
    \qquad
    Z_1,\ldots,Z_n \overset{\mathrm{iid}}{\sim} Q,
\]
with the two samples independent.
The kernel and the regularization parameter $\lambda>0$
are fixed before these samples are drawn.
\end{assumption}

For any distribution $P$, define its mean embedding and
uncentered second-moment operator by
\[
    \mu_P=\E_P[\phi(X)],
    \qquad
    T_P=\E_P[\phi(X)\otimes\phi(X)],
\]
where
\[
    (\phi(x)\otimes\phi(x))f
    =\inner{\phi(x)}{f}\phi(x).
\]
In particular,
\[
    \inner{f}{T_Pf}=\E_P[f(X)^2].
\]

For the reference sample, write
\[
    \widehat\mu_i=\frac1n\sum_{j=1}^n\phi(X_j),
    \qquad
    \widehat T_i
    =\frac1n\sum_{j=1}^n\phi(X_j)\otimes\phi(X_j).
\]
For the incoming sample, write
\[
    \widehat\mu_Q=\frac1n\sum_{j=1}^n\phi(Z_j).
\]
Set
\[
    \Delta_i=\mu_Q-\mu_{P_i},
    \qquad
    \widehat\Delta_i=\widehat\mu_Q-\widehat\mu_i,
\]
and
\[
    A_i=T_{P_i}+\lambda I,
    \qquad
    \widehat A_i=\widehat T_i+\lambda I.
\]

The population and empirical scores are
\begin{align}
    D_{\lambda,i}
    &:=\frac12\inner{\Delta_i}{A_i^{-1}\Delta_i},
    \label{eq:population-score}\\
    \widehat D_{\lambda,i}
    &:=\frac12
       \inner{\widehat\Delta_i}
             {\widehat A_i^{-1}\widehat\Delta_i}.
    \label{eq:empirical-score}
\end{align}
These quantities are well-defined because
$A_i,\widehat A_i\succeq\lambda I$.

The score is a regularized RKHS Pearson-type discrepancy.
It need not equal the classical unregularized Pearson
divergence.

For a pairwise comparison, the incoming window is accepted
as compatible with regime $i$ when
\[
    \widehat D_{\lambda,i}\le\eta,
\]
and rejected otherwise.

\section{Null guarantee and threshold}

Define, for $\rho_0\in(0,1)$,
\begin{equation}
    U_n(\rho_0)
    :=
    \frac{K}{n\lambda}
    \left(1+\sqrt{\log\frac1{\rho_0}}\right)^2.
    \label{eq:null-threshold}
\end{equation}

\begin{theorem}[Acceptance under the null]
\label{thm:null}
Under Assumption~\ref{ass:setting}, if $Q=P_i$, then
\[
    \Prob\left\{
        \widehat D_{\lambda,i}\le U_n(\rho_0)
    \right\}\ge 1-\rho_0.
\]
Consequently, any threshold $\eta\ge U_n(\rho_0)$
has pairwise false-rejection probability at most $\rho_0$.
\end{theorem}

\begin{proof}
Under the null, $\Delta_i=0$. Independence and centering give
\[
    \E\norm{\widehat\Delta_i}^2
    =
    \frac2n
    \left(
        \E_{P_i}[k(X,X)]-\norm{\mu_{P_i}}^2
    \right)
    \le\frac{2K}{n}.
\]
Hence
\[
    \E\norm{\widehat\Delta_i}
    \le\sqrt{\frac{2K}{n}}.
\]

Nonnegativity of the kernel implies
\[
    \norm{\phi(x)-\phi(y)}^2
    =
    k(x,x)+k(y,y)-2k(x,y)
    \le 2K.
\]
Replacing any one of the $2n$ observations therefore
changes $\norm{\widehat\Delta_i}$ by at most
$\sqrt{2K}/n$. The sum of squared bounded-difference
constants is $4K/n$.

One-sided McDiarmid concentration yields
\[
    \Prob\left\{
        \norm{\widehat\Delta_i}
        >
        \sqrt{\frac{2K}{n}}
        \left(1+\sqrt{\log\frac1{\rho_0}}\right)
    \right\}
    \le\rho_0.
\]
Finally, $\widehat A_i^{-1}\preceq\lambda^{-1}I$ gives
\[
    \widehat D_{\lambda,i}
    \le\frac1{2\lambda}\norm{\widehat\Delta_i}^2,
\]
which proves the result.
\end{proof}

\section{Directional guarantee under the alternative}

The alternative proof controls the empirical score along
one fixed population direction. It requires concentration
of two scalar averages rather than an operator-norm bound
for $\widehat T_i-T_{P_i}$.

\subsection{Variational representation and population direction}

For any bounded self-adjoint operator $B\succeq\lambda I$
and $v\in\HH$, completing the square gives
\begin{equation}
    \frac12\inner{v}{B^{-1}v}
    =
    \sup_{f\in\HH}
    \left\{
        \inner{f}{v}-\frac12\inner{f}{Bf}
    \right\}.
    \label{eq:variational}
\end{equation}
Indeed,
\[
    \inner{f}{v}-\frac12\inner{f}{Bf}
    =
    \frac12\inner{v}{B^{-1}v}
    -
    \frac12
    \inner{f-B^{-1}v}{B(f-B^{-1}v)}.
\]

Suppose $D_{\lambda,i}>0$ and define
\begin{equation}
    u_i
    :=
    \frac{A_i^{-1}\Delta_i}{\sqrt{2D_{\lambda,i}}}.
    \label{eq:direction}
\end{equation}
This direction depends only on the population distributions,
the kernel, and $\lambda$. It does not depend on the
observed samples.

Its normalization gives
\begin{align}
    \inner{u_i}{\Delta_i}
    &=\sqrt{2D_{\lambda,i}},
    \label{eq:direction-signal}\\
    \inner{u_i}{A_i u_i}
    &=1.
    \label{eq:direction-energy}
\end{align}
Since $A_i\succeq\lambda I$, it follows that
\begin{equation}
    \norm{u_i}^2\le\frac1\lambda.
    \label{eq:direction-norm}
\end{equation}

Define
\begin{align}
    \widehat d_i
    &:=
    \inner{u_i}{\widehat\Delta_i}
    =
    \frac1n\sum_{j=1}^n u_i(Z_j)
    -
    \frac1n\sum_{j=1}^n u_i(X_j),
    \label{eq:empirical-signal}\\
    \widehat q_i
    &:=
    \inner{u_i}{\widehat A_i u_i}
    =
    \frac1n\sum_{j=1}^n u_i(X_j)^2
    +\lambda\norm{u_i}^2.
    \label{eq:empirical-energy}
\end{align}
Their expectations satisfy
\[
    \E\widehat d_i=\sqrt{2D_{\lambda,i}},
    \qquad
    \E\widehat q_i=1.
\]

Applying \eqref{eq:variational} to the empirical score
and restricting the supremum to $f=t u_i$, $t\ge0$, gives
\begin{align}
    \widehat D_{\lambda,i}
    &\ge
    \sup_{t\ge0}
    \left\{
        t\widehat d_i-\frac{t^2}{2}\widehat q_i
    \right\}
    \notag\\
    &=
    \frac{\pos{\widehat d_i}^2}{2\widehat q_i},
    \label{eq:directional-lower-bound}
\end{align}
where $[x]_+=\max\{x,0\}$.
The denominator is strictly positive because
$u_i\ne0$ and $\widehat A_i\succeq\lambda I$.

Thus, a lower bound on the empirical score follows from
a lower bound on $\widehat d_i$ and an upper bound on
$\widehat q_i$.

\subsection{Concentration of the two scalar quantities}

Fix $\rho_1\in(0,1)$ and set
\begin{equation}
    \ell_1:=\log\frac2{\rho_1},
    \qquad
    s_n(\rho_1):=
    \sqrt{\frac{K\ell_1}{n\lambda}},
    \qquad
    r_n(\rho_1):=
    \frac K\lambda\sqrt{\frac{\ell_1}{2n}}.
    \label{eq:directional-radii}
\end{equation}

\paragraph{Directional mean difference.}
The reproducing property and the nonnegative kernel give
\[
    |u_i(x)-u_i(y)|
    \le
    \norm{u_i}\norm{\phi(x)-\phi(y)}
    \le
    \sqrt{2K}\norm{u_i}.
\]
Therefore, each of the $2n$ independent summands in
\eqref{eq:empirical-signal} has range width at most
$\sqrt{2K}\norm{u_i}/n$.

One-sided Hoeffding concentration yields, for $t>0$,
\[
    \Prob\left\{
        \widehat d_i-\E\widehat d_i\le-t
    \right\}
    \le
    \exp\left(
        -\frac{nt^2}{2K\norm{u_i}^2}
    \right).
\]
Using \eqref{eq:direction-norm}, we obtain
\begin{equation}
    \widehat d_i
    \ge
    \sqrt{2D_{\lambda,i}}
    -
    \sqrt2\,s_n(\rho_1)
    \label{eq:signal-event}
\end{equation}
with probability at least $1-\rho_1/2$.

\paragraph{Directional quadratic term.}
For every $x$,
\[
    0\le u_i(x)^2
    \le k(x,x)\norm{u_i}^2
    \le K\norm{u_i}^2.
\]
Only the empirical average in
\eqref{eq:empirical-energy} is random.
Hoeffding concentration therefore gives
\[
    \Prob\left\{
        \widehat q_i-1\ge t
    \right\}
    \le
    \exp\left(
        -\frac{2nt^2}{K^2\norm{u_i}^4}
    \right).
\]
Using \eqref{eq:direction-norm} again gives
\begin{equation}
    \widehat q_i\le1+r_n(\rho_1)
    \label{eq:energy-event}
\end{equation}
with probability at least $1-\rho_1/2$.

The two events need not be independent: they share the
reference sample. A union bound ensures that
\eqref{eq:signal-event} and \eqref{eq:energy-event}
hold simultaneously with probability at least $1-\rho_1$.

\subsection{Alternative guarantee}

\begin{assumption}[Population separation]
\label{ass:separation}
For the fixed pair $(Q,P_i)$,
\[
    D_{\lambda,i}\ge\delta>0.
\]
\end{assumption}

\begin{theorem}[Directional lower bound]
\label{thm:alternative}
Under Assumptions~\ref{ass:setting}
and~\ref{ass:separation}, define
\begin{equation}
    L_n(\delta,\rho_1)
    :=
    \frac{
        \pos{\sqrt\delta-s_n(\rho_1)}^2
    }{
        1+r_n(\rho_1)
    }.
    \label{eq:alternative-bound}
\end{equation}
Then
\[
    \Prob\left\{
        \widehat D_{\lambda,i}
        \ge L_n(\delta,\rho_1)
    \right\}
    \ge1-\rho_1.
\]
In particular, if $\eta<L_n(\delta,\rho_1)$, then
\[
    \Prob\left\{
        \widehat D_{\lambda,i}>\eta
    \right\}
    \ge1-\rho_1.
\]
\end{theorem}

\begin{proof}
On the joint event established above,
\eqref{eq:directional-lower-bound} gives
\[
    \widehat D_{\lambda,i}
    \ge
    \frac{
        \pos{
            \sqrt{2D_{\lambda,i}}
            -\sqrt2\,s_n(\rho_1)
        }^2
    }{
        2(1+r_n(\rho_1))
    }
    =
    \frac{
        \pos{\sqrt{D_{\lambda,i}}-s_n(\rho_1)}^2
    }{
        1+r_n(\rho_1)
    }.
\]
The right-hand side is nondecreasing in $D_{\lambda,i}$.
Substituting $D_{\lambda,i}\ge\delta$ proves the claim.
\end{proof}

\begin{remark}[Role of the direction]
The direction $u_i$ is used only in the proof.
The classifier still evaluates the full empirical score
\eqref{eq:empirical-score}; it does not need to know or
estimate $u_i$.

The separation assumption concerns the regularized
Pearson-type score. No additional lower bound on
$\norm{\mu_Q-\mu_{P_i}}$ is required.
\end{remark}

\section{Threshold calibration and nonempty interval}

\subsection{Separate null and alternative error levels}

The two guarantees are compatible whenever a threshold
can be chosen in
\begin{equation}
    U_n(\rho_0)
    \le\eta<L_n(\delta,\rho_1).
    \label{eq:threshold-interval}
\end{equation}

Because $U_n(\rho_0)>0$, this interval is nonempty if
and only if
\begin{equation}
    \delta>
    \left[
        s_n(\rho_1)
        +
        \sqrt{
            (1+r_n(\rho_1))U_n(\rho_0)
        }
    \right]^2.
    \label{eq:nonempty}
\end{equation}
Indeed, $U_n(\rho_0)<L_n(\delta,\rho_1)$ first requires
$\sqrt\delta>s_n(\rho_1)$, after which taking square
roots and rearranging gives \eqref{eq:nonempty}.

The choice
\[
    \eta=U_n(\rho_0)
\]
then ensures:
\begin{itemize}
    \item under $Q=P_i$, acceptance with probability
    at least $1-\rho_0$;
    \item under $D_{\lambda,i}\ge\delta$, rejection
    with probability at least $1-\rho_1$.
\end{itemize}
These are guarantees under two different hypotheses,
not a single joint probability statement.

Condition~\eqref{eq:nonempty} is equivalent to nonemptiness
of the certified interval \eqref{eq:threshold-interval}.
It is only a sufficient condition for successful
classification itself.

\subsection{Explicit formulas in the initial parameters}

Taking $\rho_0=\rho_1=\rho$, choose
\begin{equation}
    \boxed{
    \eta
    =
    \frac{K}{n\lambda}
    \left(1+\sqrt{\log\frac1\rho}\right)^2.
    }
    \label{eq:explicit-threshold}
\end{equation}

The required population separation is
\begin{equation}
    \boxed{
    \delta>
    \frac{K}{n\lambda}
    \left[
        \sqrt{\log\frac2\rho}
        +
        \left(1+\sqrt{\log\frac1\rho}\right)
        \sqrt{
            1+
            \frac K\lambda
            \sqrt{\frac{\log(2/\rho)}{2n}}
        }
    \right]^2.
    }
    \label{eq:explicit-separation}
\end{equation}

For an independently chosen threshold
$\eta\ge U_n(\rho)$, the corresponding sufficient
condition is
\begin{equation}
    \boxed{
    \delta>
    \left[
        \sqrt{\frac{K\log(2/\rho)}{n\lambda}}
        +
        \sqrt{
            \eta
            \left(
                1+\frac K\lambda
                \sqrt{\frac{\log(2/\rho)}{2n}}
            \right)
        }
    \right]^2.
    }
    \label{eq:chosen-threshold-separation}
\end{equation}

For kernels satisfying $k(x,x)=1$, take $K=1$.
For fixed $K$, $\lambda$, and $\rho$, both the threshold
in \eqref{eq:explicit-threshold} and the sufficient
separation in \eqref{eq:explicit-separation} are of
order $n^{-1}$.

\subsection{Numerical illustration}

For
\[
    n=100,\qquad K=1,\qquad
    \lambda=0.1,\qquad \rho=0.1,
\]
the quantities above are approximately
\[
    U_n(\rho)=0.63374,
    \qquad
    s_n(\rho)=0.54733,
    \qquad
    r_n(\rho)=1.22387.
\]
Thus the threshold $\eta=0.63374$ is certified whenever
\[
    \delta>3.00849.
\]

For the same $K$, $\lambda$, and $\rho$, increasing the
window length to $n=300$ gives
\[
    \eta=0.21125,
    \qquad
    \delta>0.83985.
\]
These values specify sufficient separations; they do not
establish that a particular collection of regimes
satisfies them.

\section{Scope and implementation considerations}

The results above concern a fixed pair of distributions
and fixed sampling windows.

Reference-specific kernels are permitted: for each
comparison, the same reference kernel must be used to
construct both mean embeddings and the reference
second-moment operator. Each kernel must satisfy the
stated assumptions. A kernel learned from an independent
preliminary sample can be handled by conditioning on
that sample, provided the kernel bound holds uniformly.

A bandwidth estimated from the same reference observations
used in the score requires an additional argument;
freezing it after estimation does not by itself establish
the fixed-kernel sampling assumptions.

For a memory containing several regimes, correct
assignment requires control of the relevant comparisons
simultaneously. Likewise, repeated online decisions and
windows selected by a change-point detector require
appropriate treatment of multiple comparisons, selection,
dependence, and possible regime contamination. The
pairwise results alone do not establish those guarantees.

Finally, the implementation must evaluate
\eqref{eq:empirical-score} with the same normalization
and regularization convention. Thresholds calibrated for
a different empirical quadratic score are not directly
covered by these results.

\end{document}